\song{Pro malou Lenku (Jaromír Nohavica)}
%Mara_2016_09_05

\ve{}
\ch{Ami7}Jak mi tak docházejí \ch{D7}síly, já \ch{G}pod jazykem \ch{G/F\guitarSharp}cítím \ch{Emi}síru,\\
\ch{Ami7}víru, ztrácím \ch{D7}víru, a to mě \ch{G}míchá,\\
\ch{Ami7}minomety reflektorů st\ch{D7}{ř}ílí, jsem malým \ch{G}terčem \ch{G/F\guitarSharp}na bitevním \ch{Emi}poli,\\
kdekdo mě \ch{Ami7}skolí, a to mě \ch{D7}bolí, u srdce \ch{G}píchá.

\cho{}
\ch{Ami7}Rána jsou smutnější než \ch{D7}večer, z rozbitého \ch{G}nosu \ch{G/F\guitarSharp}krev mi \ch{Emi}teče,\\
na čísle \ch{Ami7}padesátšest jedna nula \ch{D7}devět nikdo to \ch{G}nebere,\\
\ch{Ami7}rána jsou smutnější než \ch{D7}večer, na hrachu \ch{G}klečet, \ch{G/F\guitarSharp}klečet, \ch{Emi}klečet,\\
\ch{Ami7}opustil mě můj děd Vše\ch{D7}věd a zas je \ch{G}{ú}terek.\\
\ch{Ami7}Jak se ti \ch{D7}vede? No \ch{G}někdy fajn a \ch{G/F\guitarSharp}někdy je to v \ch{Emi}háji,\\
dva \ch{Ami}pozounisti z vesnické \ch{D7}kutálky pod okny mi \ch{G}hrají:\\
tú tú \ch{Ami}tú... \ch{D7}\quad{}\ch{G}\qquad{}\ch{Emi}\qquad{} \ch{Ami}\qquad{} \ch{D7}\qquad{}\ch{G}\qquad{}\ch{Ami}\qquad{}\ch{D7}\qquad{}\ch{G}\qquad{}\ch{Emi}\qquad{}\ch{Ami}\qquad{}\ch{D7}\qquad{}\ch{G} 

\ve{}
Jak říká kamarád Pepa: \uv{Co po mně chcete, slečno z první řady,\\
vaše vnady mě nebaví a trošku baví.}\\
Sudička moje byla slepá, když mi řekla to, co mi řekla,\\
píšou mi z pekla, že prý mě zdraví, že prý mě zdraví. 

\cho{} 

\ve{}
Má malá Lenko, co teď děláš, chápej, že čtyři roky, to jsou čtyři roky\\
a čas pádí a já jsem tady a ty zase jinde,\\
až umřem, říkej, žes' nás měla, to pro tebe skládáme tyhle sloky,\\
na hrachu klečíme a hloupě brečíme a světu dáváme kvinde. 

\cho{} 

\esong